\documentclass[11pt]{article}
\usepackage[a4paper,margin=25mm]{geometry}
\usepackage{amsmath,amssymb,amsthm,mathtools,booktabs,array,xcolor}
\usepackage[T1]{fontenc}\usepackage{lmodern}
\usepackage[colorlinks=true,linkcolor=teal,citecolor=teal,urlcolor=teal]{hyperref}
\hypersetup{pdftitle={Logarithmic Green Domains at Actual Whitney Cross-Caps},pdfauthor={Akari Hayami (Jian-Yu Huang)},pdfsubject={Local continuation 0.03}}
\newtheorem{theorem}{Theorem}[section]
\newtheorem{lemma}[theorem]{Lemma}
\newtheorem{corollary}[theorem]{Corollary}
\theoremstyle{remark}\newtheorem{remark}[theorem]{Remark}
\newcommand{\R}{\mathbb R}\newcommand{\C}{\mathbb C}
\newcommand{\Dom}{\operatorname{Dom}}\newcommand{\dd}{\,d}
\newcommand{\norm}[1]{\lVert#1\rVert}
\title{Logarithmic Green Domains\\at Actual Whitney Cross-Caps}
\author{Akari Hayami (Jian-Yu Huang)}
\date{2 October 2026\\\small Research continuation 0.03; local working draft}
\begin{document}\maketitle
\begin{abstract}
An arclength coordinate along the degenerate direction of a smooth Euclidean
Whitney germ makes its scalar divergence tensor equal to the identity up to
$O(\rho^{1/2})$. Subtracting the planar logarithmic fundamental solution
therefore leaves a divergence source in $L^3$ and a scalar source in $L^2$.
The scalar H\"older theorem gives a continuous regular part. For the
fixed-Dirichlet two-point restriction on the signal normalization this identifies
the actual resolvent Green vectors, proves their universal logarithmic
coefficient, and gives geometric boundary coordinates on the whole adjoint
domain. These coordinates establish Markov uniqueness and invariance under a
common change of logarithmic reference length, as well as sheet-even limiting
coefficients. The abstract extension framework is classical. The result repairs
the scalar domain statements for a specified operator; singular Dirac and Pin
transmission and publication priority remain separate questions.
\end{abstract}

\section{The operator and the comparison being made}\label{sec:scope}
The source is the completed physical rectangle of the map
\[
 S(t,u)=\bigl(c+2u(1-c),(1-u)s,u\sqrt{c(2-c)}\bigr),\qquad
 c=\cos t,\quad s=\sin t,
\]
with $-\pi/2\leq t\leq\pi/2$ and $0\leq u\leq1$. Endpoints use the
completed smooth charts of the v13 manuscript \cite{v13}. Let $P_+,P_-$ be its
two interior Whitney rank-loss points, let $g=S^*g_{\R^3}$, and let
$\mathcal H=L^2(\dd A_g)$. Our scalar product is conjugate-linear in its
first argument. The closed gradient form with fixed outer Dirichlet condition
defines the positive self-adjoint operator $H_F$. The previous continuation
\cite{v02} establishes graph-bounded, surjective values
\[
 \tau:\Dom H_F\longrightarrow\C^2,\qquad
 \tau u=(u(P_+),u(P_-)),
 \qquad \overline{\ker\tau}^{\mathcal H}=\mathcal H.
\]
Write $A=H_F|_{\ker\tau}$. It is closed, dense and symmetric, with
deficiency indices $(2,2)$. For $\eta<0$ put
\[
 R_\eta=(H_F-\eta)^{-1},\quad
 G_\eta=(\tau R_\eta)^*,\quad G_{\eta,j}=G_\eta e_j.
\]
The written Hilbert-space argument of \cite{v02}, also an instance of
Posilicano's established framework \cite{pos}, gives
\begin{equation}\label{eq:decomposition}
 \Dom A^*=\Dom H_F\dotplus G_\eta\C^2,\qquad
 A^*(u+G_\eta a)=H_Fu+\eta G_\eta a.
\end{equation}
The purpose here is to turn these resolvent coordinates into actual asymptotic
coefficients. The local argument concerns full neighborhoods of the two
interior cross-caps. No extra boundary coefficients are introduced at the
other three rank-loss labels: their Friedrichs form and the outer boundary
condition are already fixed.

This choice matters. The graph closure $H_0$ of the compactly supported
interior Laplacian leaves infinitely many outer-boundary modes free. Its
deficiency and harmonic kernel are infinite-dimensional \cite{v13}.
Statements below about $A$ cannot be read as a finite classification of $H_0$.

\section{An asymptotically Euclidean coordinate}\label{sec:coordinate}
At an actual smooth Euclidean Whitney germ, a target translation and rotation,
followed by a smooth source coordinate change, put the map in the form
\begin{equation}\label{eq:adapted}
 \begin{gathered}
 F(x,y)=(x,h(x,y)),\qquad h(0)=0,\quad Dh(0)=0,\\
 h_y(x,y)=xv+2yw+O(R^2),\qquad R=(x^2+y^2)^{1/2},
 \end{gathered}
\end{equation}
where $h$ takes values in $\R^2$ and $v,w$ are linearly independent.
Indeed, take $x$ to be the target tangent-axis projection of $F$, and choose
$y$ along its kernel at zero. The Whitney determinant
$\det(F_x,F_{xy},F_{yy})(0)\ne0$ gives the independence of $v,w$.
This uses the actual Euclidean metric; an arbitrary target diffeomorphism
would not preserve it. In a sufficiently small neighborhood,
$cR\leq|h_y|\leq CR$ off zero and $|h_x|\leq CR$.

\begin{lemma}[Arclength resolution]\label{lem:arclength}
Set
\[
 V(x,y)=\int_0^y |h_y(x,s)|\dd s,\qquad
 \Xi(x,y)=(X,V)=(x,V(x,y)),\qquad \rho=(X^2+V^2)^{1/2}.
\]
The map $\Xi$ is a local homeomorphism and a $C^1$ diffeomorphism off zero.
In its coordinates the metric matrix $G$, its density $m=\sqrt{\det G}$,
and its divergence tensor $a=mG^{-1}$ satisfy
\begin{equation}\label{eq:euclidean}
 G-I=O(\rho^{1/2}),\qquad a-I=O(\rho^{1/2}),\qquad m-1=O(\rho).
\end{equation}
They extend continuously at zero with $G(0)=a(0)=I$, $m(0)=1$, and are
bounded and uniformly elliptic near zero. The actual extrinsic radius
$\rho_e=|F(x,y)|$ satisfies
\begin{equation}\label{eq:radii}
 \rho_e/\rho=1+O(\rho^{1/2}),\qquad
 \log\rho_e-\log\rho=O(\rho^{1/2}).
\end{equation}
\end{lemma}
\begin{proof}
For fixed $x$, $V$ is strictly increasing in $y$. The bounds on $h_y$ imply
\[
 c y^2/2\leq |V(x,y)|\leq C(|xy|+y^2).
\]
Thus its inverse is continuous at zero and elsewhere, and its image contains
a neighborhood of zero. Put $n=h_y/|h_y|$ away from zero. Differentiation
under the integral, justified by the uniform bound on $h_{xy}$ and dominated
convergence, gives
\[
 V_y=|h_y|,\qquad V_x=\int_0^y n(x,s)\cdot h_{xy}(x,s)\dd s,
 \qquad |V_x|\leq C|y|.
\]
The single undefined integrand value at $(x,s)=(0,0)$ has no effect.
These derivatives are continuous, including at zero. Since $V_y>0$ off
zero, the inverse-function theorem gives the asserted $C^1$ regularity.
We need no second derivative of $V$ along $x=0$.

Let $q=h_x-V_xn$. Differentiating the actual map in the new coordinates yields
$F_X=(1,q)$, $F_V=(0,n)$, and hence
\[
 G=\begin{pmatrix}1+|q|^2&q\cdot n\\q\cdot n&1\end{pmatrix},\qquad
 \det G=1+|q|^2-(q\cdot n)^2\geq1.
\]
Here $|q|\leq CR$. Since $y^2\leq C|V|$, we have $R^2\leq C\rho$
after shrinking the neighborhood. This proves \eqref{eq:euclidean} and
uniform ellipticity, including the better density estimate.

To compare radii, Taylor's formula gives
\[
 |h-wy^2|\leq C\{x^2+|xy|+|y|R^2\},\qquad
 |V-\operatorname{sgn}(y)|w|y^2|
       \leq C\{|xy|+|y|R^2\}.
\]
The second estimate follows by integrating
$\bigl||h_y(x,s)|-2|w||s|\bigr|\leq C(|x|+x^2+s^2)$.
It follows that $\bigl||h|-|V|\bigr|\leq CR(|x|+y^2)\leq CR\rho$.
Comparing the norms of $(x,|h|)$ and $(x,|V|)$ proves
$|\rho_e-\rho|\leq CR\rho$. As $R\leq C\rho^{1/2}$, both radius
claims follow.
\end{proof}

Grieser already proved singular quasi-Euclidean coordinate resolution for the
standard Whitney metric \cite{grieser}. The present proof supplies the
asymptotic identity tensor and the extrinsic-radius normalization for an
actual Euclidean germ. This is a refinement to be compared with the literature,
not a claim to have originated singular coordinate resolution.

\section{The actual Green vectors}\label{sec:green}
We use the following external scalar regularity input \cite{simon}:
if $a$ is real, bounded and uniformly elliptic in dimension two, an $H^1$
weak solution of
$-\operatorname{div}(a\nabla w)=f_0-\operatorname{div}F_0$
is locally H\"older when $F_0\in L^q$, $f_0\in L^{q/2}$, $q>2$.
Bounded zeroth-order terms can be moved to $f_0$; here $q=3$ and $w\in L^2$.
Real and imaginary parts are treated separately. Simon's Lecture 18,
Theorems 2 and 3, printed pages 212--216, give the local boundedness and
H\"older estimates with precisely these hypotheses.

\begin{theorem}[Logarithmic normalization]\label{thm:log}
At $P_i$, with $\rho_{e,i}=|S-S(P_i)|$, there are real constants
$B_{\eta,ij}$ and $\beta>0$ such that
\begin{equation}\label{eq:greenexp}
 G_{\eta,j}=-\frac{\delta_{ij}}{2\pi}\log\rho_{e,i}
                  +B_{\eta,ij}+O(\rho_{e,i}^{\beta}).
\end{equation}
The matrix $B_\eta$ is real symmetric. The Green vector has the fixed outer
Dirichlet condition inherited from $R_\eta$.
\end{theorem}
\begin{proof}
Work first with real functions at one pole in the coordinate of
Lemma~\ref{lem:arclength}. The form is locally
$\int \nabla f\cdot a\nabla v\dd X\dd V$ and the measure is
$m\dd X\dd V$. Take a cutoff $\chi=1$ near zero, supported in this chart,
and set $h_0=-\chi\log\rho/(2\pi)$. It is in $\mathcal H$, but not in
the form domain. Distributionally $-\Delta h_0=\delta_0+k$, where $k$
is supported in the cutoff annulus and is bounded. Write $E=a-I$.
The remaining functional in
\[
 (-\operatorname{div}(a\nabla)-\eta m)h_0=\delta_0+\mathcal F
\]
is represented by $E\nabla h_0$ in the divergence term, together with $k$
and $-\eta mh_0$ in the scalar term. In particular
\[
 |E\nabla h_0|\leq C\rho^{-1/2},\qquad
 E\nabla h_0\in L^3,\qquad k-\eta mh_0\in L^2.
\]
Indeed $\int_0^\epsilon \rho^{1-3/2}\dd\rho<\infty$. The cutoff
annulus adds only bounded terms. Extending by zero makes $\mathcal F$ a
bounded functional on the global form domain. Lax--Milgram for the coercive
form $q_\eta=q-\eta\langle\cdot,\cdot\rangle$ gives a unique form-domain
$w$ with $q_\eta(w,v)=\mathcal F(v)$. The quoted scalar theorem applies
to its weak equation (including the $\eta mw$ term in $L^2$), so $w$ is
H\"older at the pole. Thus $g=h_0-w$ has a logarithm and a H\"older
regular part, and solves the point-source equation. Away from its pole it is
a weak homogeneous solution and is H\"older at the other pole as well.
Equation~\eqref{eq:radii} replaces $\log\rho$ by $\log\rho_e$, reducing
the exponent, if necessary, to $\beta\leq1/2$.

We must still identify this solution with the resolvent-defined vector;
distributional existence by itself would not suffice. For
$u\in\Dom H_F$, scalar regularity gives $|u-u(P_i)|\leq C\rho^\alpha$.
Testing its weak equation with a cutoff times $u-u(P_i)$ gives the local
Caccioppoli bound
\[
 \int_{B_r}|\nabla u|^2\leq C r^{2\alpha},\qquad 0<r<r_0,
\]
after taking $0<\alpha<1$. The gradient-cutoff term is bounded by
$Cr^{-2}\int_{B_{2r}}|u-u(P_i)|^2\leq Cr^{2\alpha}$; the scalar source
term is at most $Cr^{1+\alpha}\norm{mH_Fu}_{L^2(B_{2r})}$, which has
the same bound for small $r$. Ellipticity absorbs the cross term.

The improper integral $\int a\nabla h_0\cdot\nabla u$ therefore
converges: on dyadic annuli $\norm{\nabla\log\rho}_{L^2}$ is bounded,
while $\norm{\nabla u}_{L^2}=O(r^\alpha)$, a summable bound. Let
$\theta_\epsilon$ vanish on $B_\epsilon$ and equal one off $B_{2\epsilon}$,
with $|\nabla\theta_\epsilon|\leq C/\epsilon$.
The form test $h_0\theta_\epsilon$ in the equation for $u$ has an extra
term bounded by $C|\log\epsilon|\epsilon^\alpha\to0$.
Consequently
\[
 \langle h_0,(H_F-\eta)u\rangle
       =u(P_i)+\mathcal F(u).
\]
For the identity on the right, the planar radial flux of
$-\log\rho/(2\pi)$ has total mass one; replacing $u$ by its circle mean
gives $u(P_i)$ in the limit, with error $O(\epsilon^\alpha)$.
The $E$ and scalar terms converge in their form-dual pairings. The equation
for $w$ similarly gives
$\langle w,(H_F-\eta)u\rangle=\mathcal F(u)$.
Thus $\langle g,(H_F-\eta)u\rangle=u(P_i)$ for every $u\in\Dom H_F$.
Taking $u=R_\eta f$ identifies $g=(\tau_iR_\eta)^*1=G_{\eta,i}$.
Complex functions follow by linearity with the stated scalar-product convention.

Reality of the form and $\eta$ makes all $G_{\eta,j}$ and their regular
constants real. For reciprocity, let $\delta_{\epsilon,i}$ be the normalized
indicator of a resolved ball about $P_i$, normalized using $\dd A_g$.
Graph H\"older estimates imply
$\langle\delta_{\epsilon,i},\cdot\rangle\to\tau_i$ in the dual graph norm,
so $R_\eta\delta_{\epsilon,i}\to G_{\eta,i}$ in $\mathcal H$.
Near a different pole these functions solve the homogeneous equation;
the local H\"older and boundedness estimates improve convergence to uniform
convergence on a smaller ball. Self-adjointness of $R_\eta$ gives
\[
 \langle\delta_{\epsilon,i},R_\eta\delta_{\epsilon,j}\rangle
 =\langle\delta_{\epsilon,j},R_\eta\delta_{\epsilon,i}\rangle.
\]
For $i\ne j$, passage to the limit gives
$G_{\eta,j}(P_i)=G_{\eta,i}(P_j)$. The diagonal entries are already real,
so $B_\eta$ is real symmetric.
\end{proof}

\section{Geometric coordinates on the whole adjoint domain}\label{sec:domains}
\begin{theorem}[Expansion and Green pairing]\label{thm:boundary}
Every $f\in\Dom A^*$ has unique coefficients $\ell,b\in\C^2$ with
\begin{equation}\label{eq:maxexp}
 f=\ell_i\log\rho_{e,i}+b_i+O(\rho_{e,i}^{\beta})\quad\hbox{at }P_i.
\end{equation}
For $f=u+G_\eta a$ in \eqref{eq:decomposition}, they are
\[
 \ell=-a/(2\pi),\qquad b=\tau u+B_\eta a.
\]
Both coefficient maps are bounded in the graph norm of $A^*$, and the combined
map is onto $\C^2\oplus\C^2$, with kernel $\Dom A$. For $f,g$ with
coefficients $(\ell,b)$ and $(\ell',b')$,
\begin{equation}\label{eq:greenpair}
 \langle A^*f,g\rangle-\langle f,A^*g\rangle
 =2\pi\{\langle\ell,b'\rangle-\langle b,\ell'\rangle\}.
\end{equation}
The Friedrichs domain is precisely $\ell=0$. All self-adjoint extensions are
the inverse images of maximal isotropic planes for \eqref{eq:greenpair}.
Equivalently they have, for a unique $U\in U(2)$,
\[
 (U-I)b+i(U+I)a=0,\qquad a=-2\pi\ell.
\]
Here $U=I$ gives $H_F$. This parametrization differs by the regular-part
shift $B_\eta$ from the resolvent coordinates in \cite{v02}.
\end{theorem}
\begin{proof}
Theorem~\ref{thm:log}, the graph H\"older estimate for $u$, and
\eqref{eq:decomposition} give the expansion. Its uniqueness follows by
first dividing by $\log\rho_e$ and then taking the remaining limit.
Moreover $u=R_\eta(A^*-\eta)f$ is graph-bounded, and the finite-dimensional
inverse on $G_\eta\C^2$ bounds $a$. Surjectivity follows by prescribing $a$
and then using surjectivity of $\tau$ to prescribe $b-B_\eta a$.
Their common kernel is $a=0$, $\tau u=0$, namely $\Dom A$.

The actual resolvent decomposition gives the pairing
$\langle\tau u,a'\rangle-\langle a,\tau u'\rangle$.
Substituting $\tau u=b-B_\eta a$ cancels the $B_\eta$ terms by symmetry
and gives \eqref{eq:greenpair}. Vanishing $a$ is exactly
$\Dom H_F$ in the direct sum. The finite-dimensional maximal-isotropic
classification now applies to this surjective actual boundary map; its Cayley
description gives the displayed $U(2)$ equation.
\end{proof}

For $z\in\rho(H_F)$, retain the previous Weyl matrix
$M_z=(z-\eta)\tau R_zG_\eta$. The actual geometric regular-part matrix is
\begin{equation}\label{eq:regularmatrix}
 B_z=B_\eta+M_z,
\end{equation}
because $G_z-G_\eta=(z-\eta)R_zG_\eta\in\Dom H_F$.
If a self-adjoint matrix $C$ specifies $b=Ca$, the resolvent is
\[
 (H_C-z)^{-1}=R_z+G_z(C-B_z)^{-1}G_{\bar z}^*,
\]
whenever the inverse exists. The Krein extension of $A$ has plane $b=B_0a$
and kernel $G_0\C^2$. These are consequences of the existing extension
framework, with a now-proved interpretation of its coefficients.

\begin{corollary}[Actual reflection symmetry]\label{cor:reflection}
For real $\eta<0$ the matrix $B_\eta$ has the form
$\left(\begin{smallmatrix}r_\eta&g_\eta\\g_\eta&r_\eta\end{smallmatrix}\right)$.
The same holds at $z=0$.
\end{corollary}
\begin{proof}
The actual map satisfies $S(-t,u)=J S(t,u)$, where
$J(X,Y,Z)=(X,-Y,Z)$ is Euclidean orthogonal. Source reflection preserves
the induced form, area and fixed outer Dirichlet condition, so its unitary
pullback commutes with $H_F$ and $R_\eta$ and exchanges $G_{\eta,+}$ and
$G_{\eta,-}$. It exchanges the extrinsic radii exactly. The two diagonal
constants are consequently equal; reciprocity supplies the equal real
off-diagonal constants. The argument also applies to $R_0$ since $H_F$ has
a positive gap. This symmetry exchanges the two poles; it is distinct from
the local sheet exchange at a single cross-cap.
\end{proof}

\section{Two uniqueness statements and local parity}\label{sec:uniqueness}
\begin{corollary}[Markov uniqueness for the specified restriction]\label{cor:markov}
$H_F$ is the only self-adjoint extension of $A$ whose heat semigroup is
submarkovian.
\end{corollary}
\begin{proof}
The closed gradient form with outer Dirichlet condition is a Dirichlet form:
normal contractions decrease the gradient energy and preserve its closure.
Thus its operator $H_F$ generates a submarkovian semigroup.
Conversely, for any such self-adjoint extension $H$, $\lambda>0$, and
$h\in L^2\cap L^\infty$, the Laplace-transform formula for its resolvent
gives
\[
 \norm{(H+\lambda)^{-1}h}_\infty\leq\lambda^{-1}\norm h_\infty.
\]
By \eqref{eq:maxexp}, any nonzero logarithmic coefficient makes this
resolvent vector essentially unbounded near its pole; the regular remainder
is bounded and the resolved measure has a positive bounded density.
Hence every coefficient vanishes. The vector lies in $\Dom H_F$, and
$(H+\lambda)^{-1}h=(H_F+\lambda)^{-1}h$. Bounded $L^2$ inputs are dense,
so the resolvents, and therefore the self-adjoint operators, are equal.
\end{proof}

\begin{corollary}[Common reference-length invariance]\label{cor:scale}
Among self-adjoint boundary planes for $A$, the only one invariant under all
simultaneous changes of logarithmic reference length is the Friedrichs plane.
\end{corollary}
\begin{proof}
Writing $f=\ell_i\log(\rho_{e,i}/\lambda_i)+b_i^{(\lambda)}+o(1)$ gives
$b^{(\lambda)}=b+\operatorname{diag}(\log\lambda_i)\ell$.
The change $\lambda_i\mapsto e^t\lambda_i$ acts by
$(\ell,b)\mapsto(\ell,b+t\ell)$. If a maximal isotropic plane $L$ is
invariant for every real $t$, then $(\ell,b)\in L$ implies $(0,\ell)\in L$.
Pairing these vectors in \eqref{eq:greenpair} yields
$2\pi\norm\ell^2=0$, so $\ell=0$. Dimension two gives the entire
Friedrichs plane. It is itself invariant.
\end{proof}
This statement changes a reference convention. It does not assert a dilation
isometry of the actual surface or a physical renormalization law.

\begin{corollary}[Sheet-even limiting coefficients]\label{cor:parity}
If paired source points $p_\epsilon,q_\epsilon$ approach the same interior
cross-cap and $S(p_\epsilon)=S(q_\epsilon)$, then every
$f\in\Dom A^*$ satisfies
$f(p_\epsilon)-f(q_\epsilon)=O(\rho_e^\beta)\to0$.
Thus the logarithmic and regular boundary coefficients are sheet-even.
\end{corollary}
\begin{proof}
The common image has exactly the same extrinsic radius on both branches.
Their logarithmic terms in \eqref{eq:maxexp} cancel, their constant terms
coincide, and both remainders tend to zero with the stated bound.
\end{proof}
This is asymptotic parity of scalar coefficients. It requires no sheet
isometry of the actual metric and asserts no invariance of the full Green
function. The ordinary support obstruction and the odd-to-even equivariant
linear obstruction of v13 remain applicable. No canonical derived-sheaf-to-
operator morphism, phase-fiber map or singular Pin transmission is supplied.

\section{What has been repaired, and what is still missing}\label{sec:comparison}
The ancestral TeX and the authoritative final v12 are different records.
The 1,908-line ancestor supplies the earlier signal motivation; the 43-page
final v12 adds the late scalar and spinorial assertions. This continuation
proves logarithmic coefficients and scalar uniqueness for $A$, replacing the
ill-defined finite classification of the old $H_0$. It recovers the intended
two-point conclusion under a specified fixed outer boundary, not under the
old compact-core definition.

The construction in continuation 0.02 uses component powers $(1-u,u)$.
The ancestral moment assignment was $(\cos t,1-\cos t)$. They are different
state families. The explicit new quadrature realization proves that $S$ can
be realized with supplied phases; it does not derive $S$ from the ancestral
moment map alone. The Hermitian no-go also explains why changing this
readout into three pure-state quantum expectations is impossible.

Abstract boundary triples and resolvent formulas are established theory
\cite{pos}. Markov uniqueness for degenerate operators also has a substantial
earlier literature \cite{robsik}; its theorem for a compact core on an open
Euclidean domain is not imported here, because our outer boundary is fixed
and its coefficient hypotheses differ. The potential contribution lies in
the actual Euclidean Whitney estimates, coefficient normalization and the
fully specified application. Independent derivation and a new proof in this
project do not establish publication priority. That question needs further
comparison with singular-metric and point-interaction literature.

The principal mathematical tasks still open are the actual singular Dirac
graph domains, conformal and gauge estimates, front/seam overlap control,
and a proved phase-fiber/Pin transmission map. The scalar H\"older theorem
does not prove a Dirac system estimate. The earlier weight-zero commutator
problem and the missing fiber identification remain unresolved. Physical
clock, spin, chirality, channel capacity or arithmetic realizations need
their own objects and assumptions. The new scalar PDE and operator proofs
are written proofs using external theorems; the existing 33-theorem Lean
inventory has not been enlarged or reinterpreted as full formalization.

\begin{thebibliography}{9}
\bibitem{v13} A. Hayami (J.-Y. Huang), \emph{Bicomplex Signal Geometry:
 Completed Source Charts, Monodromy and Spectral Stability},
 v13 local working manuscript, 1 October 2026. Preserved companion 0.01.
\bibitem{v02} A. Hayami (J.-Y. Huang), \emph{Bicomplex Signal Realization
 and Fixed-Boundary Point Extensions}, local continuation 0.02, 1 October 2026.
\bibitem{grieser} D. Grieser, \emph{Quasiisometry of singular metrics},
 Houston J. Math. \textbf{28} (2002), 741--752.
 \url{https://www.math.uh.edu/~hjm/Vol28-4.html}.
\bibitem{pos} A. Posilicano, \emph{Boundary Triples and Weyl Functions for
 Singular Perturbations of Self-Adjoint Operators}, arXiv preprint 2003,
 revised 15 January 2004.
 \url{https://arxiv.org/abs/math/0309077}. Sections 2--3, Theorems 2.2 and 3.1.
\bibitem{simon} L. Simon, \emph{Lectures on PDE}, author notes, updated
 5 March 2015, Lecture 18, Theorems 2--3, printed pp.212--216.
 \url{https://math.stanford.edu/~lms/lecs-on-pde.pdf}.
\bibitem{robsik} D. W. Robinson and A. Sikora, \emph{Markov uniqueness of
 degenerate elliptic operators}, Ann. Sc. Norm. Super. Pisa Cl. Sci.
 (5) \textbf{10} (2011), 683--710.
 \url{https://arxiv.org/abs/0912.4536}. Related theory; not used as a theorem
 for the resolved coefficients or the fixed-Dirichlet restriction.
\end{thebibliography}
\end{document}
